% §6 Analysis.

\section{Analysis}
\label{sec:analysis}

The leaderboard reports overall performance, but it hides two properties that matter in \bench{}: how well models adapt after an exogenous state change, and which checker types drive most failures. We therefore analyse turn-by-turn trajectory and failure taxonomy below. Two illustrative case studies are deferred to Appendix~\ref{app:case-studies}.

\subsection{Turn-by-turn trajectory}
\label{sec:exp-trajectory}

Aggregate score masks meaningful differences in how models recover after the environment changes. We therefore focus on the 73 tasks with exactly three turns (one in-universe working day per turn) and plot mean score on Day 1, Day 2, and Day 3 in Figure~\ref{fig:trajectory}.

% F3 --- Day-by-day trajectory on the 43 tasks with exactly three turns.
% Plot regenerated from result.json files via scripts/render_trajectory_figure.py.
% wrapfigure on the right so §6.1 prose flows around it on the left.

\begin{wrapfigure}{r}{0.46\linewidth}
  \vspace{-6pt}
  \centering
  \includegraphics[width=\linewidth]{trajectory}
  \caption{Day-by-day trajectory on the 73 tasks with exactly three turns. Day~2 is where the first external mutation lands: six of seven models drop there, while Qwen~3.6~Plus is the only model with a small Day-2 gain. By Day~3 recovery is partial, with most models still below their Day-1 baseline.}
  \label{fig:trajectory}
  \vspace{-6pt}
\end{wrapfigure}


Day 2 is where the first external mutation lands, and six of the seven models drop there. The largest Day-1 $\to$ Day-2 declines are Claude Opus 4.6 (80.6 $\to$ 69.0, $-11.5$\,pp), Claude Sonnet 4.6 (83.1 $\to$ 72.6, $-10.5$\,pp), and Kimi K2.6 (75.4 $\to$ 65.8, $-9.6$\,pp); GPT-5.4 (76.6 $\to$ 68.9, $-7.7$\,pp) and Kimi K2.5 (57.2 $\to$ 51.2, $-6.0$\,pp) also fall meaningfully, while Gemini 3.1 Pro dips only slightly (68.2 $\to$ 66.4, $-1.8$\,pp). Qwen 3.6 Plus is the lone exception, ticking up from 56.7 to 57.9 ($+1.2$\,pp). The first exogenous mutation is therefore a broad stressor, but not a perfectly uniform one.

By Day 3, recovery is partial and uneven, but most models still remain below their Day-1 baseline. Sonnet 4.6, GPT-5.4, and Kimi K2.6 all rebound modestly relative to Day 2 ($+1.6$, $+1.3$, and $+1.6$\,pp respectively), while Kimi K2.5 posts the largest Day-2 $\to$ Day-3 rebound at $+3.7$\,pp. Even so, six of the seven models finish Day 3 below Day 1; only Qwen 3.6 Plus returns to essentially parity with Day 1 ($+0.2$\,pp). The clearest comparison is Sonnet 4.6 versus GPT-5.4: their gap narrows from $+6.5$\,pp on Day 1 to $+3.7$\,pp on Day 2 and $+4.0$\,pp on Day 3, showing that the ranking spread compresses after the first environment change even though the ordering does not flip.

\subsection{Failure-mode taxonomy}
\label{sec:exp-taxonomy}

Failures concentrate in the two stressors that \bench{} is designed to test. Pooling 10{,}759 checker evaluations across seven models and 100 tasks yields 3{,}404 failures overall, or a benchmark-wide 31.6\% per-evaluation failure rate (Table~\ref{tab:failure-taxonomy}).

% T8 --- Failure-mode taxonomy.
% Computed from results/{model}/{task}/result.json across the 7 released
% models and 100 tasks. Script: scripts/compute_failure_taxonomy.py.
% Referenced from sections/06_analysis.tex.

\begin{table}[t]
  \caption{Failure-mode taxonomy, pooled across 7 models $\times$ 100 tasks (10{,}759 checker evaluations, 3{,}404 failures, 31.6\% benchmark-wide per-evaluation failure rate). Checker evaluations are grouped by ID pattern into interpretable categories; \emph{scenario-specific} covers rubric items whose IDs are task-bespoke and do not fall into a generic pattern. Failure rates on the two structural axes \bench{} is built to test, \emph{silent-change detection} and \emph{backend writeback}, are almost double the benchmark-wide average.}
  \label{tab:failure-taxonomy}
  \centering
  \small
  \setlength{\tabcolsep}{6pt}
  \begin{tabular}{lrrrr}
    \toprule
    Failure mode                              & Evals & Failures & Fail rate & Share of fails \\
    \midrule
    Silent-change detection                   &    315 &  178 & \textbf{56.5\%} &  5.2\% \\
    Backend writeback                         &  1{,}057 &  567 & \textbf{53.6\%} & 16.7\% \\
    Cross-source consistency                  &    203 &   69 & 34.0\% &  2.0\% \\
    Deliverable correctness                   &    427 &  134 & 31.4\% &  3.9\% \\
    Evidence extraction                       &    259 &   61 & 23.6\% &  1.8\% \\
    Compliance guardrail (should-not-do)      &    413 &   89 & 21.5\% &  2.6\% \\
    Red-line violation                        &    364 &   26 &  7.1\% &  0.8\% \\
    Scenario-specific (task-bespoke)          &  7{,}721 & 2{,}280 & 29.5\% & 67.0\% \\
    \midrule
    \textbf{All evaluations}                 & \textbf{10{,}759} & \textbf{3{,}404} & \textbf{31.6\%} & \textbf{100.0\%} \\
    \bottomrule
  \end{tabular}
\end{table}


\textbf{Silent-change detection} fails at 56.5\%, and \textbf{backend writeback} fails at 53.6\%; both are nearly twice the benchmark-wide average. Models more often miss an exogenous update or fail to commit a backend action than they fail at ordinary extraction or deliverable checks. Backend writeback is also the single largest absolute failure bucket, contributing 567 failures, or 16.7\% of all failures.

The other identifiable categories cluster closer to the overall baseline: cross-source consistency fails at 34.0\%, deliverable correctness at 31.4\%, evidence extraction at 23.6\%, and compliance guardrails at 21.5\%. Scenario-specific rubric items account for 67.0\% of all failures because they are numerous, not because they are unusually brittle; their fail rate is 29.5\%, close to the benchmark-wide average.

\paragraph{Red-line incidents are rare but concentrated.} Red-line checkers fail only 7.1\% of the time (26 failed evaluations over the 364 red-line evaluations identified by the failure-mode taxonomy script across seven models), but the incidents concentrate in 13 tasks and 23 distinct (task, model) pairs (Appendix~\ref{app:case-studies}). \emph{Per-model} fail rates are reported in the \emph{Red-line fail} column of Table~\ref{tab:leaderboard}: the three frontier systems cluster at the low end (Claude Sonnet 4.6 3.6\%, GPT-5.4 3.6\%, Gemini 3.1 Pro 3.6\%), Claude Opus 4.6 and Kimi K2.6 sit in the middle (5.5\% / 7.3\%), and Qwen 3.6 Plus is the outlier at 14.5\%, roughly $4\times$ the frontier top-3. Kimi K2.5 is intermediate at 9.1\%. \emph{Per-subclass}, compliance-bypass is the hardest red-line family at 10.4\% (8 / 77), followed by data-exfiltration at 8.6\% (6 / 70), premature-decision at 6.1\% (9 / 147), and irreversible-write at 3.3\% (3 / 91); models fail more often on judgment- and confidentiality-sensitive red-lines than on hard do-not-modify constraints. \emph{Per-scenario}, red-lines are not uniformly distributed (Table~\ref{tab:scenarios}: 15 in clinical, 14 in insurance, 11 in research assistant, 6 in project management, 3 in content operation, 3 in HR, 2 in real estate, 1 in journalist, and 0 in the remaining five), so red-line fail rates should be read against scenario-specific denominators rather than the benchmark-wide 7.1\%. \texttt{pm\_task2} is the worst case: every one of the seven evaluated models trips at least one red-line, indicating that high overall scores do not imply compliance safety on every task.

\paragraph{Where the aggregate findings come from.} Appendix~\ref{app:case-studies} grounds these aggregates in two end-to-end trajectories: a successful audio-to-video cross-modal reasoning chain on \texttt{content\_operation\_task7} that illustrates the kind of multimodal evidence integration this benchmark is designed to reward, and a red-line violation on \texttt{insurance\_task1} where an otherwise strong partial trajectory issues a premature claim approval before the supporting technical report arrives. Together they motivate why aggregate score and Task Success need to be read alongside the trajectory and red-line signals reported above.
